\documentclass{beamer}\usepackage[]{graphicx}\usepackage[]{color}
%% maxwidth is the original width if it is less than linewidth
%% otherwise use linewidth (to make sure the graphics do not exceed the margin)
\makeatletter
\def\maxwidth{ %
  \ifdim\Gin@nat@width>\linewidth
    \linewidth
  \else
    \Gin@nat@width
  \fi
}
\makeatother

\definecolor{fgcolor}{rgb}{0.345, 0.345, 0.345}
\newcommand{\hlnum}[1]{\textcolor[rgb]{0.686,0.059,0.569}{#1}}%
\newcommand{\hlstr}[1]{\textcolor[rgb]{0.192,0.494,0.8}{#1}}%
\newcommand{\hlcom}[1]{\textcolor[rgb]{0.678,0.584,0.686}{\textit{#1}}}%
\newcommand{\hlopt}[1]{\textcolor[rgb]{0,0,0}{#1}}%
\newcommand{\hlstd}[1]{\textcolor[rgb]{0.345,0.345,0.345}{#1}}%
\newcommand{\hlkwa}[1]{\textcolor[rgb]{0.161,0.373,0.58}{\textbf{#1}}}%
\newcommand{\hlkwb}[1]{\textcolor[rgb]{0.69,0.353,0.396}{#1}}%
\newcommand{\hlkwc}[1]{\textcolor[rgb]{0.333,0.667,0.333}{#1}}%
\newcommand{\hlkwd}[1]{\textcolor[rgb]{0.737,0.353,0.396}{\textbf{#1}}}%

\usepackage{framed}
\makeatletter
\newenvironment{kframe}{%
 \def\at@end@of@kframe{}%
 \ifinner\ifhmode%
  \def\at@end@of@kframe{\end{minipage}}%
  \begin{minipage}{\columnwidth}%
 \fi\fi%
 \def\FrameCommand##1{\hskip\@totalleftmargin \hskip-\fboxsep
 \colorbox{shadecolor}{##1}\hskip-\fboxsep
     % There is no \\@totalrightmargin, so:
     \hskip-\linewidth \hskip-\@totalleftmargin \hskip\columnwidth}%
 \MakeFramed {\advance\hsize-\width
   \@totalleftmargin\z@ \linewidth\hsize
   \@setminipage}}%
 {\par\unskip\endMakeFramed%
 \at@end@of@kframe}
\makeatother

\definecolor{shadecolor}{rgb}{.97, .97, .97}
\definecolor{messagecolor}{rgb}{0, 0, 0}
\definecolor{warningcolor}{rgb}{1, 0, 1}
\definecolor{errorcolor}{rgb}{1, 0, 0}
\newenvironment{knitrout}{}{} % an empty environment to be redefined in TeX

\usepackage{alltt} 
% \usepackage{graphicx}
\usepackage{graphics}
\usepackage[T1]{fontenc}

\setbeamercovered{transparent}


\renewcommand{\ni}{\noindent}


% \SweaveOpts{cache=TRUE, background="white"}


\title[ 1-Intro]{1 - Intro to R}
\subtitle{Data Structures}
\author[ K. Maurer, S. VanderPlas, E. Hare]{Karsten Maurer, Eric Hare, and Susan VanderPlas}
\date{February 12, 2014}
\institute[ISU]{Iowa State University}
\IfFileExists{upquote.sty}{\usepackage{upquote}}{}

\begin{document}





\begin{frame}
    \maketitle
\end{frame}

%--------------------------------------------------------------------------------------

\begin{frame}
    \frametitle{Data Frames}
    \begin{itemize}
        \item Data Frames are the work horse of R objects
        \item Structured by rows and columns and can be indexed
        \item Each column is a specified variable type
        \item Columns names can be used to index a variable
        \item Advice for naming variable applys to editing columns names
        \item Can be specified by grouping vectors of equal length as columns
    \end{itemize}    
\end{frame}

\begin{frame}[fragile]
    \frametitle{Data Frame Indexing}
    \begin{itemize}
        \item Elements indexed similar to a vector using \texttt{[} \texttt{]}
        \item \texttt{df[i,j]} will select the element in the $i^{th}$ row and $j^{th}$ column
        \item \texttt{df[ ,j]} will select the entire $j^{th}$ column and treat it as a vector
        \item \texttt{df[i ,]} will select the entire $i^{th}$ row and treat it as a vector
        \item Logical vectors can be used in place of i and j used to subset the row and columns
    \end{itemize}
\end{frame}

%--------------------------------------------------------------------------------------

\begin{frame}
    \frametitle{Adding a New Variable to Data Frames}
    \begin{itemize}
        \item Create a new vector that is the same length as other columns
        \item Append new column to the data frame using the \texttt{\$} operator
        \item The new data frame column will adopt the name of the vector
    \end{itemize}
\end{frame}


%--------------------------------------------------------------------------------------

\begin{frame}[fragile]
    \frametitle{Data Frame Demo}
    \begin{itemize}
        \item Demo using a statistical classic: Edgar Anderson's Iris Data
        \item Follow along with the R script named 4-DataStructures.R
    \end{itemize}
\end{frame}

%--------------------------------------------------------------------------------------

\begin{frame}
    \frametitle{Your Turn}
    \begin{itemize}
        \item Make a data frame with column 1: 1,2,3,4,5,6 and column 2:a,b,a,b,a,b
        \item Select only rows with value "a" in column 2 using logical vector
        \item \texttt{mtcars} is a built in data set like \texttt{iris}: read the values from row 4 by indexing
    \end{itemize}    
\end{frame}

%--------------------------------------------------------------------------------------

\begin{frame}
    \frametitle{Lists}
    \begin{itemize}
        \item Lists are a structured collection of R objects
        \item R objects in a list need not be the same type
        \item Create lists using the \texttt{list()} function
        \item Lists indexed using double square brackets [[ ]] to select an object
    \end{itemize}    
\end{frame}

%--------------------------------------------------------------------------------------

\begin{frame}[fragile]
    \frametitle{List Example}
    \small
\begin{knitrout}\footnotesize
\definecolor{shadecolor}{rgb}{1, 1, 1}\color{fgcolor}\begin{kframe}
\begin{alltt}
\hlcom{# Create a list including a vector and a matrix}
\hlstd{mylist} \hlkwb{<-} \hlkwd{list}\hlstd{(}\hlkwd{matrix}\hlstd{(letters[}\hlnum{1}\hlopt{:}\hlnum{10}\hlstd{],} \hlkwc{nrow} \hlstd{=} \hlnum{2}\hlstd{,} \hlkwc{ncol} \hlstd{=} \hlnum{5}\hlstd{),} \hlkwd{seq}\hlstd{(}\hlnum{0}\hlstd{,} \hlnum{49}\hlstd{,} \hlkwc{by} \hlstd{=} \hlnum{7}\hlstd{))}
\hlstd{mylist}
\end{alltt}
\begin{verbatim}
## [[1]]
##      [,1] [,2] [,3] [,4] [,5]
## [1,] "a"  "c"  "e"  "g"  "i" 
## [2,] "b"  "d"  "f"  "h"  "j" 
## 
## [[2]]
## [1]  0  7 14 21 28 35 42 49
\end{verbatim}
\begin{alltt}
\hlcom{# Use index to select second object in list, this will return the vector}
\hlstd{mylist[[}\hlnum{2}\hlstd{]]}
\end{alltt}
\begin{verbatim}
## [1]  0  7 14 21 28 35 42 49
\end{verbatim}
\end{kframe}
\end{knitrout}

    \normalsize
\end{frame}

%--------------------------------------------------------------------------------------

\begin{frame}
    \frametitle{Your Turn}
    \begin{itemize}
        \item Create a list containing a vector and a 2x3 data frame
        \item Use indexing to select the data frame from your list
        \item Use further indexing to select the first row from the data frame in your list
    \end{itemize}    
\end{frame}

%--------------------------------------------------------------------------------------

\begin{frame}
    \frametitle{Examining Objects}
    \begin{itemize}
        \item \texttt{head(x)} - View top 6 rows of a data frame
        \item \texttt{tail(x)} - View bottom 6 rows of a data frame
        \item \texttt{summary(x)} - Summary statistics   
        \item \texttt{str(x)} - View structure of object  
        \item \texttt{dim(x)} - View Dimensions of object 
        \item \texttt{length(x)} - Returns the length of a vector
    \end{itemize}    
\end{frame}

%--------------------------------------------------------------------------------------

\begin{frame}[fragile]
    \frametitle{Examining Objects Example}
    \footnotesize
\begin{knitrout}\footnotesize
\definecolor{shadecolor}{rgb}{1, 1, 1}\color{fgcolor}\begin{kframe}
\begin{alltt}
\hlcom{# Examine the top 2 rows of the iris data set from built in data package}
\hlkwd{head}\hlstd{(iris,} \hlnum{2}\hlstd{)}
\end{alltt}
\begin{verbatim}
##   Sepal.Length Sepal.Width Petal.Length Petal.Width Species
## 1          5.1         3.5          1.4         0.2  setosa
## 2          4.9         3.0          1.4         0.2  setosa
\end{verbatim}
\begin{alltt}
\hlcom{# How big is this data set?}
\hlkwd{dim}\hlstd{(iris)}
\end{alltt}
\begin{verbatim}
## [1] 150   5
\end{verbatim}
\begin{alltt}
\hlcom{# What structure does the data set have?}
\hlkwd{str}\hlstd{(iris)}
\end{alltt}
\begin{verbatim}
## 'data.frame':	150 obs. of  5 variables:
##  $ Sepal.Length: num  5.1 4.9 4.7 4.6 5 5.4 4.6 5 4.4 4.9 ...
##  $ Sepal.Width : num  3.5 3 3.2 3.1 3.6 3.9 3.4 3.4 2.9 3.1 ...
##  $ Petal.Length: num  1.4 1.4 1.3 1.5 1.4 1.7 1.4 1.5 1.4 1.5 ...
##  $ Petal.Width : num  0.2 0.2 0.2 0.2 0.2 0.4 0.3 0.2 0.2 0.1 ...
##  $ Species     : Factor w/ 3 levels "setosa","versicolor",..: 1 1 1 1 1 1 1 1 1 1 ...
\end{verbatim}
\end{kframe}
\end{knitrout}

    \normalsize
\end{frame}

%--------------------------------------------------------------------------------------

\begin{frame}
    \frametitle{Your Turn}
    \begin{itemize}
        \item View the top 6 rows of mtcars data
        \item What type of object is the mtcars data set?
        \item How many rows are in iris data set? (try finding this using dim or indexing + length)
        \item Summarize the values in each column in iris data set
    \end{itemize}    
\end{frame}

%--------------------------------------------------------------------------------------

\begin{frame}
    \frametitle{Working with output from a function}
    \begin{itemize}
        \item Can save output from a function as an object
        \item Object is generally a list of output objects
        \item Can pull off items from the output for further computing
        \item Examine object using functions like \texttt{str(x)}
    \end{itemize}
\end{frame}

%--------------------------------------------------------------------------------------

\begin{frame}
    \frametitle{Output Object Demo}
    \begin{itemize}
        \item Demo of saving t-test output as an object
    \end{itemize}
\end{frame}

%--------------------------------------------------------------------------------------

\begin{frame}
    \frametitle{Your Turn}
    \begin{itemize}
        \item Pull the p-value from the t-test of a difference between Sepal Lengths of setosa and versicolor species from the Iris data
    \end{itemize}
\end{frame}

%--------------------------------------------------------------------------------------

\begin{frame}
    \frametitle{Importing Data}
    \begin{itemize}
        \item First need to tell R where the data is saved using \texttt{setwd()}
        \item Data read in using R functions such as:
            \begin{itemize}
            \item \texttt{read.table()} for reading in .txt files
            \item \texttt{read.csv()} for reading in .csv files
            \end{itemize}
        \item Assign the data to new R object when reading in the file
    \end{itemize}
\end{frame}

%--------------------------------------------------------------------------------------

\begin{frame}
    \frametitle{Importing Data Demo}
    \begin{itemize}
        \item Demo of creating a csv file and loading it into R
    \end{itemize}
\end{frame}

%--------------------------------------------------------------------------------------

\begin{frame}
    \frametitle{Your Turn}
    \begin{itemize}
        \item Make 5 rows of data in an excel spreadsheet and save it as a .txt file
        \item Import this new .txt file into R with read.table
    \end{itemize}    
\end{frame}

\end{document}
